% Pillar-5 (MIN-REPORT-RL) evidence exhibits
\section{Evidence: The Stack Moves More Than the Label}
\label{sec:p5-evidence}

\paragraph{Versioning note.}
The frozen manifest corpus predates the evaluation-report addition to
MIN-REPORT-RL.  Accordingly, the audits below score the standard's first seven
\emph{run-manifest fields}; the current standard adds pass@$k$ curves as
Item~8 (\secref{sec:p5-stack}).  Exhibits~18--20 also evaluate a historical
continuous-telemetry extension.  That experiment is reported as an optional
extension and must not be read as redefining the current Item~8.

\paragraph{Evidence boundary.}
Only Exhibits 1--3 below motivate this reporting proposal. The later historical
mega-corpus, badge-extension, and cross-corpus variance-ratio analyses are not
retained as validation. The 98-cell mega corpus is training-free sampling,
not an RL algorithm comparison. Marginal one-way eta-squared values are not
additive causal variance partitions, and ratios across different experimental
populations do not rank stack versus algorithm importance.

\subsection{Exhibit 1: A \texorpdfstring{17$\times$}{17x} Same-Label Gap Is Mostly a Checkpoint Switch}
\label{sec:p5-exhibit-backend}
The framework-comparison record (\texttt{experiments/results/framework\_comparison.json};
Stratum B) labels two runs ``GRPO, Qwen3-8B, GSM8K, $G{=}8$, lr $10^{-5}$,
30 steps, seed 42''. Final (last-10) training reward is $0.856$ on Tinker and
$0.050$ on TRL: a $17\times$ ratio. The runs are not a backend comparison:
\begin{itemize}
  \item \textbf{Checkpoint.} The Tinker run trained \texttt{Qwen3-8B-Base};
    the TRL run trained the Instruct \texttt{Qwen3-8B}. A run of Qwen3-8B (Instruct) on the same
    Tinker harness
    (\texttt{wave6\_w6\_base\_r32\_t0.8\_b2} in \texttt{master\_results.json})
    reaches $0.338$. The checkpoint switch alone therefore moves reward by
    $-0.519$, i.e.\ $64\%$ of the $-0.806$ gap; the residual $-0.288$
    (Tinker vs.\ TRL, both on Qwen3-8B (Instruct)) is the most that could be
    attributed to ``backend''.
  \item \textbf{Everything else that still differs.} Reward function and
    answer parser; system prompt and chat template; decoding (temperature,
    maximum completion length); batch shape ($2$ prompts $\times 8$ on Tinker
    vs.\ $1\times8$ on TRL); training-data subset (TRL used
    \texttt{train[:500]}); and loss/optimizer implementation (LoRA rank and
    targets, clipping, normalization). The TRL run logs reward $0.0$ on its
    first step and on 25 of its 30 steps, consistent with a reward/format
    mismatch rather than a learning-dynamics difference.
  \item \textbf{The ``matched'' configuration.} The per-framework
    ``matched-configuration'' YAML in \texttt{experiments/framework\_config\_dumps/}
    describes neither run: it was written post hoc for reviewers (with
    placeholder checkpoint hashes), and, e.g., it lists LoRA rank 16 and 384
    new tokens for TRL whereas the TRL script used rank 32 and 512.
\end{itemize}
An earlier revision stated that every labelled knob except the checkpoint was
harmonised. That is false and is withdrawn. The exhibit supports one narrow
point: a label plus a short hyperparameter table did not identify the model or
the harness actually compared, and the largest single lever was the
unreported checkpoint. A designed decomposition (same checkpoint, reward,
prompt, decoding, and batch across both trainers) has not been run.

\subsection{Exhibit 2: The Same Label on Different Experiments}
\label{sec:p5-exhibit-label}
DAPO~\citep{yu2025dapo} is defined by decoupled (asymmetric) clipping
\emph{plus} dynamic sampling, which resamples all-correct and all-incorrect
groups and so drives the Zero-Variance Fraction (ZVF) to zero by construction.
An earlier revision contrasted an open-trainer ``DAPO'' run (mean ZVF $0.00$)
with a closed-stack ``DAPO'' run (mean ZVF $0.58$) as ``same label, opposite
telemetry''. That contrast is not interpretable: the two runs differ in
model, task (synthetic arithmetic vs.\ GSM8K), $G$, number of steps, and
seeds; the open-trainer $0.00$ holds by construction; and the closed-stack
``dapo'' arm is an advantage-clipping variant of GRPO with no dynamic
sampling, i.e.\ not DAPO, and its registry entry should not carry the DAPO
label. The defensible content is \emph{within-stack}: switching
GRPO$\to$``DAPO'' moved mean ZVF from $0.25$ to $0.00$ on the open trainer
and from $0.567$ to $0.578$ on the closed stack. The lesson for reporting is
that a label must be accompanied by the mechanism actually implemented
(here: whether dynamic sampling ran).

\subsection{Exhibit 3: Four Labels on One Fixed Stack}
\label{sec:p5-exhibit-h2h}
In a 12-cell head-to-head on the Tinker stack (Qwen3.5-4B, GSM8K, three seeds
per arm, 15 steps, batch 2, $G{=}8$; Stratum B), mean last-10 training reward
spans $0.034$ ($0.710$--$0.744$; \tableref{tab:p5-h2h}). An earlier revision
called this ``within noise''. The data do not show that. Treating arms as
independent, the pooled within-arm seed SD is $0.105$ and two one-sided
equivalence tests (TOST, $\alpha=0.05$) fail for every pair at margins
$\pm0.05$, $\pm0.10$, and $\pm0.15$; equivalence of all six pairs is reached
only at $\pm0.19$, wider than any effect of interest. Because each seed fixes
the prompt stream, the appropriate analysis blocks on seed: seed accounts for
$96\%$ of the sum of squares and arm for $2.5\%$ (arm $F_{3,6}=2.62$,
$p=0.15$); the residual SD is $0.017$; and two seed-paired contrasts exclude
zero (dapo$-$gspo $+0.033$, 90\% CI $[+0.006,+0.060]$; drgrpo$-$gspo
$+0.031$, $[+0.004,+0.058]$). All-pairs equivalence under the block model
holds only at $\pm0.060$. With 3 seeds and 15 steps, label differences of a
few points can neither be excluded nor confirmed. Mean ZVF ($0.500$--$0.578$)
is far from the $0.00$ that true dynamic sampling yields, consistent with
the ``dapo'' arm not running dynamic sampling (Exhibit~2).

\begin{table}[t]
\centering
\caption{12-cell head-to-head on one fixed stack (Tinker; Qwen3.5-4B; GSM8K;
3 seeds per arm; 15 steps; internal program records, June 2026). Differences
are \emph{not} shown to be within noise: see text for unpaired TOST and
seed-blocked contrasts.}
\label{tab:p5-h2h}
\small
\begin{tabular}{@{}lcc@{}}
\toprule
Arm (label as configured) & Mean last-10 reward & Mean ZVF \\
\midrule
dapo (advantage-clip variant; no dynamic sampling) & 0.744 & 0.578 \\
drgrpo                           & 0.742 & 0.500 \\
grpo                             & 0.723 & 0.567 \\
gspo                             & 0.710 & 0.511 \\
\bottomrule
\end{tabular}
\end{table}

\endinput % Exhibits 4--21 are withdrawn pending estimand/provenance reconciliation.